\section{Introduction}\label{sec:introduction}

The next operating system may not manage files.  It may manage meaning.

Every previous generation of operating systems arose in response to a
specific bottleneck.  As that bottleneck shifted, the entire computing
stack reorganised around a new abstraction layer.  We argue that the
same shift is happening now---and that its consequences are more
profound than either of the two that preceded it.

%% ────────────────────────────────────────────────────────────────────
\subsection{The Three Ages of Computing}\label{sec:three-ages}

Computing history divides naturally into three eras, each defined by a
dominant bottleneck, a characteristic abstraction, and a corresponding
operating-system paradigm~\cite{silberschatz2018operating}.

\begin{figure}[htbp]
\centering
\begin{tikzpicture}[
    era/.style={
      rectangle, rounded corners=6pt, draw=neosblue!60, fill=neosblue!6,
      minimum width=4.2cm, minimum height=2.6cm, align=center,
      font=\small,
    },
    label/.style={font=\footnotesize\itshape, text=gray!70!black},
    >=Stealth
  ]

  %% --- timeline ---
  \draw[gray!50, very thick] (-6.5,0) -- (6.5,0);
  \foreach \x/\yr in {-4.5/1950s, 0/1980s, 4.5/2025} {
    \fill[neosblue!50] (\x,0) circle (3pt);
    \node[below=4pt, font=\footnotesize\bfseries] at (\x,0) {\yr};
  }

  %% --- era boxes ---
  \node[era] (hw) at (-4.5,2.4)
    {\textbf{Hardware Age}\\[3pt]
     Bottleneck: \emph{the machine}\\
     OS shares CPU \& memory};

  \node[era] (sw) at (0,2.4)
    {\textbf{Software Age}\\[3pt]
     Bottleneck: \emph{complexity}\\
     OS organises programs};

  \node[era, draw=neosblue!80, fill=neosblue!12] (ia) at (4.5,2.4)
    {\textbf{Intelligence Age}\\[3pt]
     Bottleneck: \emph{meaning}\\
     OS organises thinking};

  %% --- interface labels ---
  \node[label, below=6pt] at (hw.south) {Punch cards, CLI};
  \node[label, below=6pt] at (sw.south) {GUI, IDEs, Apps};
  \node[label, below=6pt] at (ia.south) {Natural language, NEOS Shell};

  %% --- arrows ---
  \draw[->, thick, gray!60] (hw.east) -- (sw.west);
  \draw[->, thick, gray!60] (sw.east) -- (ia.west);
\end{tikzpicture}
\caption{The three ages of computing.  Each era is defined by a
  dominant bottleneck and a corresponding OS paradigm.  The Intelligence
  Age shifts the bottleneck from hardware or software to meaning
  itself.}
\label{fig:three-ages}
\end{figure}

In the \textbf{Hardware Age} (1950s--1980s), the CPU was the scarcest
resource.  Operating systems were born to answer a single question:
\emph{How do we share this machine?}  Every cycle was a resource to be
scheduled; every byte was territory to be managed.  Batch processing,
time-sharing, and virtual memory were all answers to the hardware
bottleneck.

In the \textbf{Software Age} (1980s--2020s), hardware became abundant,
but complexity became the enemy.  C, Java, Python---humans learned
programming languages to organise millions of lines of code.  Operating
systems evolved to answer: \emph{How do we organise these programs?}
Graphical user interfaces, multitasking, and application ecosystems
arose to tame software complexity.

Now the bottleneck has moved again.  The challenge is no longer
computing faster or organising more code---it is organising
\textbf{thinking itself}.  How do you represent an idea so that it can
interact with other ideas?  How do you track which interpretations are
gaining strength?  How do you merge two lines of reasoning?  How do you
\emph{debug a thought}?  These are the questions that an OS for the
Intelligence Age must answer.  They are precisely the questions NEOS was
designed to address.

\begin{calloutbox}
The central thesis of this paper is that the next abstraction layer
will not manage code---it will manage \textbf{reasoning} as a
first-class computational primitive.
\end{calloutbox}

%% ────────────────────────────────────────────────────────────────────
\subsection{The Vision: LLM as Universal OS}\label{sec:vision}

Today, large language models are \emph{applications} that run on an
operating system---Windows, Linux, macOS~\cite{brown2020language}.
They are apps.  Useful apps, even transformative apps.  But still:
apps.

\textbf{Tomorrow, the LLM \emph{is} the operating system.}  Every
interaction with a computer will pass through language understanding.
This is not speculation---it is the trajectory we are already on.  And
NEOS is the first specification for what that world looks like.

\begin{figure}[htbp]
\centering
\begin{tikzpicture}[
    ibox/.style={
      rectangle, rounded corners=6pt, draw=#1!60, fill=#1!8,
      minimum width=2.2cm, minimum height=0.7cm,
      align=center, font=\small
    },
    ibox/.default=gray,
    arr/.style={-Stealth, thick, neosblue!40},
  ]
  %% --- Interface evolution ---
  \node[ibox=gray]   (cli)  at (0,0)   {\textbf{CLI}\\[-1pt]\tiny 1970s};
  \node[ibox=gray]   (gui)  at (3,0)   {\textbf{GUI}\\[-1pt]\tiny 1984};
  \node[ibox=gray]   (nui)  at (6,0)   {\textbf{NUI}\\[-1pt]\tiny 2007};
  \node[rectangle, rounded corners=6pt, draw=neosblue!60, fill=neosblue!8, minimum width=2.2cm, minimum height=0.7cm, align=center, font=\small]   (neos) at (9.5,0) {\textbf{NEOS Shell}\\[-1pt]\tiny 2025};

  \draw[arr] (cli)  -- (gui);
  \draw[arr] (gui)  -- (nui);
  \draw[arr] (nui)  -- (neos);

  \node[below=3pt, font=\tiny\itshape, text=gray] at (cli.south)
    {Text commands};
  \node[below=3pt, font=\tiny\itshape, text=gray] at (gui.south)
    {Point \& click};
  \node[below=3pt, font=\tiny\itshape, text=gray] at (nui.south)
    {Touch \& gesture};
  \node[below=3pt, font=\tiny\itshape, text=gray] at (neos.south)
    {Meaning \& resonance};
\end{tikzpicture}
\caption{Interface evolution from command-line to meaning-based
  interaction.  Each transition raises the level of abstraction at which
  humans communicate with machines.}
\label{fig:interface-evolution}
\end{figure}

Consider what this inversion means in concrete terms:

\begin{center}
\small
\begin{tabular}{@{}ll@{}}
\toprule
\thead{Old World} & \thead{New World} \\
\midrule
File systems organise by path
  & Semantic fields organise by meaning \\
Compiled binaries execute instructions
  & Field configurations execute reasoning \\
Separate applications for each task
  & Task definitions compose into workflows \\
GUIs mediate human--machine interaction
  & Natural language + shell for precision \\
\bottomrule
\end{tabular}
\end{center}

\begin{figure}[htbp]
\centering
\begin{tikzpicture}[
    layer/.style={
      rectangle, rounded corners=6pt, draw=#1!60, fill=#1!8,
      minimum width=5cm, minimum height=1.1cm,
      align=center, font=\small
    },
    layer/.default=gray,
    lbl/.style={font=\small\bfseries, text=gray!60!black},
    >=Stealth
  ]

  %% ===== TODAY stack (left) =====
  \node[lbl] at (-3.5,4.6) {TODAY};

  \node[layer=gray] (thw)  at (-3.5,0)   {Hardware};
  \node[layer=gray] (tos)  at (-3.5,1.4) {Operating System};
  \node[layer=gray] (tapp) at (-3.5,2.8) {Applications};

  %% Small LLM box inside Applications
  \node[rectangle, rounded corners=2pt, draw=neosblue!50, fill=neosblue!10,
        minimum width=1.6cm, minimum height=0.45cm,
        font=\tiny\bfseries, text=neosblue!70!black]
    (llmbox) at (-3.5,2.8) {LLM App};

  %% ===== TOMORROW stack (right) =====
  \node[lbl] at (3.5,4.6) {TOMORROW};

  \node[rectangle, rounded corners=6pt, draw=neosblue!50, fill=neosblue!8,
        minimum width=5cm, minimum height=1.1cm, align=center, font=\small]
    (nfound) at (3.5,0) {\textbf{LLM} (Foundation Model)};
  \node[rectangle, rounded corners=6pt, draw=neosblue!60, fill=neosblue!8,
        minimum width=5cm, minimum height=1.1cm, align=center, font=\small] (nneos) at (3.5,1.4)
    {\textbf{NEOS} (Operating System)};
  \node[rectangle, rounded corners=6pt, draw=neosgreen!50!black, fill=neosgreen!5,
        minimum width=5cm, minimum height=1.1cm, align=center, font=\small]
    (nhum) at (3.5,2.8) {\textbf{Human} (Operator)};

  %% --- Arrow between stacks ---
  \draw[->, ultra thick, neosblue!30] (0.3,1.4) -- (0.8,1.4);

  %% ===== REPL Cycle below =====
  \node[font=\footnotesize\bfseries, text=neosblue!60!black]
    at (0, -1.5) {The NEOS REPL Cycle};

  \node[rectangle, rounded corners=6pt, draw=neosblue!40, fill=neosblue!4,
        minimum width=2.1cm, minimum height=0.65cm,
        font=\scriptsize, align=center] (r1) at (-4.2,-2.8) {Inject\\Meaning};
  \node[rectangle, rounded corners=6pt, draw=neosblue!40, fill=neosblue!4,
        minimum width=2.1cm, minimum height=0.65cm,
        font=\scriptsize, align=center] (r2) at (-1.4,-2.8) {Run\\Dynamics};
  \node[rectangle, rounded corners=6pt, draw=neosblue!40, fill=neosblue!4,
        minimum width=2.1cm, minimum height=0.65cm,
        font=\scriptsize, align=center] (r3) at (1.4,-2.8) {Observe\\Emergence};
  \node[rectangle, rounded corners=6pt, draw=neosblue!40, fill=neosblue!4,
        minimum width=2.1cm, minimum height=0.65cm,
        font=\scriptsize, align=center] (r4) at (4.2,-2.8) {Collapse to\\Insight};

  \draw[->, thick, neosblue!50] (r1) -- (r2);
  \draw[->, thick, neosblue!50] (r2) -- (r3);
  \draw[->, thick, neosblue!50] (r3) -- (r4);
  \draw[->, thick, neosblue!50, rounded corners=8pt]
    (r4.south) -- ++(0,-0.55) -| (r1.south);
\end{tikzpicture}
\caption{The architecture inversion.  \textbf{Left:} today's stack,
  where the LLM is an application running on a traditional OS.
  \textbf{Right:} the NEOS stack, where the LLM \emph{is} the
  substrate and NEOS provides the operating-system layer.
  \textbf{Bottom:} the NEOS REPL cycle---inject meaning, run dynamics,
  observe emergence, collapse to insight---with a feedback loop through
  save, fork, and share.}
\label{fig:architecture-inversion}
\end{figure}

The NEOS shell---\texttt{nf}---is the \textbf{REPL of cognitive
computing}.  The operator does not write code; instead, they inject
meaning, run dynamics, observe emergence, and collapse to insight
(Figure~\ref{fig:architecture-inversion}).  Chain-of-thought
reasoning~\cite{wei2022chain} demonstrated that LLMs can decompose
complex problems step by step; NEOS generalises this insight into a
full dynamical framework where reasoning steps are not linear chains
but interacting fields that self-organise through resonance.

``Programs'' in NEOS are not instruction sequences.  They are
\textbf{field configurations}---patterns with activation strengths,
resonance relationships, and tuneable parameters.  A program is a
thought, formalised.
